% Source: https://artofproblemsolving.com/community/c3h1360282_trapezoid_problem
% Stable identity: aops-topic-1360282-q01
Through $D$ draw the line parallel to $BC$, meeting $AB$ at $P$. Then $PBCD$ is a parallelogram, so $PB=CD=6$, $DP=CB=15$ and $AP=20-6=14$. Triangle $APD$ has sides $13$, $14$, $15$ and the same height $h$ as the trapezoid.

\begin{asy}
usepackage("amsmath");
size(7cm);
defaultpen(fontsize(10pt));
path ra(pair P,pair u,pair v,real s=0.2){return P+s*unit(u)--P+s*unit(u)+s*unit(v)--P+s*unit(v);}
pair A=(0,0), B=(20,0), C=(11,12), D=(5,12), P=(14,0), H=(5,0);
draw(A--B--C--D--cycle,linewidth(0.9));
draw(D--P,red+dashed); draw(D--H,blue+dashed); draw(ra(H,(1,0),(0,1),0.9),blue);
dot(P); dot(H);
label("$A$",A,SW); label("$B$",B,SE); label("$C$",C,NE); label("$D$",D,NW); label("$P$",P,S); label("$H$",H,S);
label("$6$",(C+D)/2,N); label("$13$",(A+D)/2,W); label("$15$",(B+C)/2,E);
draw((0,-1.6)--(14,-1.6),Arrows(5)); draw((14,-1.6)--(20,-1.6),Arrows(5));
draw((0,-0.6)--(0,-2.1),gray); draw((14,-1.3)--(14,-1.9),gray); draw((20,-0.6)--(20,-2.1),gray);
label("$14$",(7,-1.6),S,fontsize(9pt)); label("$6$",(17,-1.6),S,fontsize(9pt));
label("$15$",(D+P)/2,E,red); label("$12$",(D+H)/2,W,blue);
label("\mbox{}",(min(currentpicture,user=true)+max(currentpicture,user=true))/2,black); // keeps the text after the figure black when LaTeX typesets the labels inline
\end{asy}

By Heron's formula with $s=\dfrac{13+14+15}{2}=21$, the area of triangle $APD$ is $\sqrt{21\cdot8\cdot7\cdot6}=84$, so
\[
h=\dfrac{2\cdot84}{14}=12 .
\]
The trapezoid's area is $\dfrac{AB+CD}{2}\cdot h=\dfrac{20+6}{2}\cdot12=\boxed{156}$.

In general: the legs and the difference of the bases form a triangle with the same height as the trapezoid; find that height (Heron's formula, or the Pythagorean theorem with the two horizontal offsets), then use $\dfrac{a+b}{2}h$. (Here the foot $H$ of the height from $D$ has $AH=5$ and $HP=9$.)
